\documentclass[a4paper]{article}

\newcommand{\notetitle}{Lecture 01: Introducción a los sistemas de frontera — Context, Compute y la gran transformación de la infraestructura}
\newcommand{\noteauthors}{Elaborado a partir de la clase impartida por Anjney Midha, el video oficial de Stanford Online y una transcripción manual}
\newcommand{\notedate}{Primavera de 2026}
\newcommand{\videochannel}{Stanford Online / Anjney Midha}
\newcommand{\videopublishdate}{2026-04-30}
\newcommand{\videoduration}{1:05:55}
\newcommand{\videourl}{https://www.youtube.com/watch?v=O5PfU_uDhS0}
\newcommand{\videocoverpath}{cover.jpg}

% Espanol (Mexico) con XeLaTeX: fontspec para fuentes Unicode, polyglossia
% para la division silabica y los nombres del documento en la variante mexicana.
\usepackage{fontspec}
\usepackage{polyglossia}
\setdefaultlanguage[variant=mexican]{spanish}
% polyglossia no traduce los nombres de listings.
\AtBeginDocument{\providecommand{\lstlistingname}{}\renewcommand{\lstlistingname}{Listado}%
  \providecommand{\lstlistlistingname}{}\renewcommand{\lstlistlistingname}{Índice de listados}}
\usepackage{amsmath,amssymb}
\usepackage{graphicx}
\usepackage[margin=2.35cm]{geometry}
\usepackage[most]{tcolorbox}
\usepackage{etoolbox}
\usepackage{listings}
\usepackage{xcolor}
\usepackage{hyperref}
\usepackage{booktabs,longtable,tabularx,array}
\usepackage{subcaption}
\usepackage{float}
\usepackage{enumitem}
\usepackage{microtype}
\usepackage{fancyhdr}
\usepackage{needspace}

\definecolor{kbBack}{HTML}{F2F6FA}
\definecolor{kbFrame}{HTML}{9DB4CC}
\definecolor{kbTitle}{HTML}{52708F}
\definecolor{ibBack}{HTML}{FBF5E7}
\definecolor{ibFrame}{HTML}{C9A961}
\definecolor{ibTitle}{HTML}{7D5F1E}
\definecolor{wbBack}{HTML}{FCF2F0}
\definecolor{wbFrame}{HTML}{D6A096}
\definecolor{wbTitle}{HTML}{9E4F43}
\definecolor{dbBack}{HTML}{F4F7F3}
\definecolor{dbFrame}{HTML}{AEC2AC}
\definecolor{dbTitle}{HTML}{5F7F64}

\tcbset{softbox/.style={
  enhanced,breakable,boxrule=0.8pt,arc=2pt,left=3mm,right=3mm,
  top=2mm,bottom=2mm,coltitle=white,fonttitle=\bfseries,
  toptitle=0.6mm,bottomtitle=0.6mm
}}
\newtcolorbox{knowledgebox}[1]{softbox,colback=kbBack,colframe=kbFrame,colbacktitle=kbTitle,title={#1}}
\newtcolorbox{importantbox}[1]{softbox,colback=ibBack,colframe=ibFrame,colbacktitle=ibTitle,title={#1}}
\newtcolorbox{warningbox}[1]{softbox,colback=wbBack,colframe=wbFrame,colbacktitle=wbTitle,title={#1}}
\newtcolorbox{dialoguebox}[1]{softbox,colback=dbBack,colframe=dbFrame,colbacktitle=dbTitle,title={#1}}

\lstset{
  language=Python,basicstyle=\ttfamily\small,
  keywordstyle=\color{kbTitle}\bfseries,stringstyle=\color{wbTitle},
  commentstyle=\color{dbTitle},breaklines=true,frame=single,numbers=left,
  numberstyle=\tiny\color{gray},captionpos=b,columns=fullflexible,
  keepspaces=true,showstringspaces=false,extendedchars=true
}
\BeforeBeginEnvironment{lstlisting}{\Needspace{12\baselineskip}}

\hypersetup{colorlinks=true,linkcolor=kbTitle,urlcolor=kbTitle,citecolor=wbTitle}
\setlist{itemsep=0.25em,topsep=0.4em}
\setlength{\parindent}{2em}
\setlength{\parskip}{0.25em}
\newcolumntype{L}[1]{>{\raggedright\arraybackslash}p{#1}}

\providecommand{\notetitle}{Notas del curso: [tema del curso]}
\providecommand{\noteauthors}{Elaboradas a partir de los materiales públicos de Stanford CS153}
\providecommand{\notedate}{Primavera de 2026}
\providecommand{\videochannel}{Stanford Online}
\providecommand{\videopublishdate}{2026}
\providecommand{\videoduration}{[duración del video]}
\providecommand{\videourl}{}
\providecommand{\videocoverpath}{cover.jpg}
\providecommand{\coursesite}{https://cs153.stanford.edu/}
\providecommand{\repourl}{https://github.com/hqhq1025/ai-course-notes}
\providecommand{\skillversion}{wdkns/wdkns-skills@39f1a04}

\newcommand{\figsource}[1]{\par\vspace{-0.3em}{\footnotesize\color{gray}\emph{Fuente: #1}}\par}
\newcommand{\teachervoice}[1]{\begin{knowledgebox}{Nota de clase}#1\end{knowledgebox}}
\newcommand{\term}[1]{\textbf{#1}}

\pagestyle{fancy}
\fancyhf{}
\fancyfoot[C]{\thepage}
\fancyfoot[R]{\footnotesize\color{gray}\href{\repourl}{AI Course Notes}}
\renewcommand{\headrulewidth}{0pt}
\renewcommand{\footrulewidth}{0pt}

\newcommand{\makecscover}{
\begin{titlepage}
\centering
\vspace*{0.7cm}
{\Large Stanford CS153 · Frontier Systems\par}
\vspace{0.8cm}
{\huge\bfseries \notetitle\par}
\vspace{0.6cm}
{\large \noteauthors\par}
\vspace{0.25cm}
{\large \notedate\par}
\vspace{0.9cm}
\ifdefempty{\videocoverpath}{}{
  \includegraphics[width=0.86\textwidth,height=0.43\textheight,keepaspectratio]{\videocoverpath}\par
}
\vfill
\begin{tcolorbox}[width=0.92\textwidth,colback=black!2!white,colframe=black!45,boxrule=0.7pt,arc=2pt]
\textbf{Autor o canal del video}: \videochannel\par
\textbf{Fecha de publicación}: \videopublishdate\par
\textbf{Duración del video}: \videoduration\par
\textbf{Sitio del curso}: \href{\coursesite}{\nolinkurl{\coursesite}}\par
\ifdefempty{\videourl}{}{\textbf{Enlace del video}: \href{\videourl}{\nolinkurl{\videourl}}\par}
\textbf{Normas de elaboración}: \skillversion\ + las normas source-first / teacher-voice / visual-QA de este repositorio
\end{tcolorbox}
\end{titlepage}
\tableofcontents
\newpage
}

\graphicspath{{./}{images/}}
\setcounter{tocdepth}{2}

\begin{document}
\makecscover

\section{El curso no es un directorio de celebridades: primero, un método para estudiarlo}

Esta clase inaugural no empieza por los modelos de GPU, el entrenamiento distribuido ni las API de cómputo en la nube, sino por una pregunta más básica: ¿cómo debe estudiar un alumno un curso en el que participa una gran cantidad de líderes de la industria? El ponente se presenta como el opening act de toda la temporada, llama headliners a los invitados que vendrán después y asume por iniciativa propia la broma de ``AI Coachella''. Esta apertura recuerda al lector que el elenco por sí mismo no es conocimiento; lo que de verdad aporta valor es comparar cómo los distintos actores definen los problemas, cómo exponen sus restricciones y cómo su posición de intereses influye en sus juicios.

\begin{figure}[H]
\centering
\includegraphics[width=0.9\textwidth]{images/00-02-45-life-scaling-laws.jpg}
\caption{El ponente toma prestado el lenguaje de las scaling laws para ofrecer heurísticas de vida sobre las relaciones, el disfrute y la influencia a largo plazo.}
\figsource{Video oficial de Stanford Online 00:02:45, `images/00-02-45-life-scaling-laws.jpg`.}
\end{figure}

\begin{knowledgebox}{Lectura de la figura: por qué un curso técnico empieza hablando de ``life scaling laws''}
Esta diapositiva no disfraza la vida de una curva de entrenamiento que pueda predecirse con precisión. El ponente subraya que sus reglas empíricas provienen de la observación personal: dedicarse en serio al trabajo, conservar el disfrute y construir relaciones con personas con quienes uno quiera convivir a largo plazo. La palabra clave es \term{empirical}, es decir, válido según la experiencia, pero sin garantía de cumplirse para cada persona ni en cada etapa. Dejar clara esta limitación prepara, de hecho, la discusión posterior sobre el scaling de la IA: una curva puede describir las regularidades ya observadas, pero no demuestra por sí sola que el futuro vaya a prolongarse indefinidamente.
\end{knowledgebox}

\begin{warningbox}{No confundir una metáfora con una ley}
``La vida es como una tarea de entrenamiento'' solo sirve para ordenar las variables de decisión; no puede darle a la vida un objective único que se pueda optimizar. La salud, las relaciones, las responsabilidades y la influencia no comparten una misma escala. El ponente conserva la sencilla regla empírica de ``have fun with people you enjoy'' justamente para evitar que todas las decisiones se reduzcan a cifras de currículum, salario o influencia.
\end{warningbox}

\subsection{Preparedness: qué busca formar el curso}

Esta sección cambia el objetivo de aprendizaje de «escuchar las opiniones de figuras célebres» a la \term{preparedness}, es decir, el grado de preparación para enfrentar sistemas reales y restricciones organizacionales y sociales. Estar preparado no equivale a memorizar las conclusiones de los invitados; se acerca más a tres capacidades: descomponer una afirmación ambiciosa en supuestos técnicos y económicos verificables; ver las dependencias que están fuera del nivel en el que uno trabaja; y mantener la incertidumbre cuando la evidencia es insuficiente. Por ello, las clases con invitados de las diez semanas siguientes deben tomarse como un conjunto de casos en vivo con posturas distintas, no como un conjunto de respuestas unificadas.

\begin{figure}[H]
\centering
\includegraphics[width=0.86\textwidth]{images/00-07-45-introduction-to-anj.jpg}
\caption{La trayectoria académica y profesional de Anjney Midha, y su participación en las etapas iniciales de varias organizaciones de IA.}
\figsource{Video oficial de Stanford Online 00:07:45, `images/00-07-45-introduction-to-anj.jpg`.}
\end{figure}

Esta diapositiva debe leerse como una declaración sobre el origen de los puntos de vista. La trayectoria del ponente abarca formación en matemáticas, ciencias computacionales, economía y bioinformática en Stanford, e incluye organizaciones como Discord, Anthropic, Mistral, Black Forest Labs, AMP y Periodic Labs. Este historial lo vuelve más sensible a la manera en que los resultados de investigación se convierten en empresas, productos e infraestructura, pero también implica que sus juicios provienen de una perspectiva de inversión, emprendimiento y despliegue. El lector debe asimilar la experiencia práctica y, al mismo tiempo, mantener una auditoría del sesgo de muestra y de la posición de intereses.

\begin{knowledgebox}{Nota de clase: el objetivo del curso no es optimizar la búsqueda de prácticas profesionales}
El ponente distingue explícitamente entre internship y preparedness. La primera tiende a reducir cada clase a «qué perfil contrata esta empresa»; la segunda exige entender por qué una organización invierte en determinado nivel, qué recursos de las etapas anteriores la restringen y qué costos traslada a las etapas posteriores. Por eso, al hacer el proyecto del curso, la definición del problema, la cadena de evidencia y un prototipo verificable importan más que perseguir términos de moda.
\end{knowledgebox}

\begin{importantbox}{Marco de lectura de esta clase}
Todos los casos que siguen se auditan con cinco preguntas: quién controla los recursos escasos; cómo entran esos recursos en el ciclo cerrado del producto; de qué retroalimentación depende la expansión; qué puede demostrar la evidencia disponible; y quién coordina los estándares, los mercados o las instituciones. Este marco sirve tanto para el entrenamiento de modelos como para el mercado de GPU, los servicios en la nube y la infraestructura pública.
\end{importantbox}

\subsection{Resumen de la sección}

La experiencia de vida de la apertura, la trayectoria del ponente y el elenco de invitados definen en conjunto el método del curso: valorar la práctica de primera mano sin idolatrar la identidad de nadie; usar regularidades empíricas sin disfrazarlas de predicciones seguras; y ampliar el objetivo de aprendizaje de las oportunidades a corto plazo al juicio sobre sistemas que atraviesan varios niveles. A continuación se aborda el verdadero eje técnico de esta clase: por qué la IA está reabriendo una pila de sistemas que alguna vez fue relativamente estable.
\section{Del stack de sistemas estable a la Great Transition}

Durante más de una década, el software de internet consolidó una vía de producción relativamente estable: el capital adquiere terrenos, electricidad y centros de datos; los chips llegan a las salas de servidores; las plataformas en la nube empaquetan el hardware como recursos invocables; los equipos entrenan modelos y construyen aplicaciones sobre ellas, y después los mecanismos de seguridad, cumplimiento normativo y gobernanza llevan el sistema al mundo real. La IA generativa no se limitó a añadir una categoría de aplicaciones en la cima de este stack; al mismo tiempo elevó la presión sobre las capas de consumo energético, interconexión, chips, datos, modelos, productos y gobernanza, lo que obliga a renegociar los antiguos límites de responsabilidad y estructuras de precios.

\begin{figure}[H]
\centering
\includegraphics[width=0.94\textwidth]{images/00-10-00-course-in-a-nutshell.jpg}
\caption{El curso organiza el contenido de diez semanas con una cadena de todo el stack que va del capital, la energía y los centros de datos a los modelos, las aplicaciones y la gobernanza.}
\figsource{Video oficial de Stanford Online 00:10:00, `images/00-10-00-course-in-a-nutshell.jpg`.}
\end{figure}

\begin{knowledgebox}{Lectura de la figura: primero en la dirección de las dependencias, después en la dirección del valor}
De abajo hacia arriba, el capital, los terrenos, la electricidad, los chips y los centros de datos aportan las condiciones físicas para el cómputo; la cloud/software infrastructure convierte el hardware heterogéneo en recursos orquestables; los modelos y los agent consumen esos recursos para producir capacidades; las aplicaciones incorporan esas capacidades a los flujos de trabajo, y la capa de gobernanza se ocupa de la seguridad, la responsabilidad y la aceptación social. De arriba hacia abajo, los ingresos de las aplicaciones y el valor público determinan, a su vez, en qué tipo de capacidad de cómputo, modelos e instalaciones está dispuesto el capital a seguir invirtiendo. Ambas direcciones coexisten, por lo que cualquier «avance» en una sola capa puede quedar frenado por otra.
\end{knowledgebox}

\begin{table}[H]
\centering
\begin{tabular}{L{0.18\textwidth}L{0.34\textwidth}L{0.38\textwidth}}
\toprule
Capa & Recursos o interfaces principales & Preguntas de auditoría de Frontier Systems \\
\midrule
Capital y energía & Financiamiento, terrenos, red eléctrica, compra de electricidad a largo plazo & ¿El ciclo de inversión puede ajustarse a la incertidumbre de la tecnología y de la demanda? \\
Centros de datos y chips & Salas de servidores, aceleradores, redes, enfriamiento & ¿Qué hardware es realmente intercambiable y dónde hay dependencia de proveedor? \\
Nube e infraestructura de software & Planificación, virtualización, almacenamiento, observabilidad & ¿Cómo se convierten los recursos heterogéneos en interfaces de servicio estables? \\
Modelos y agent & Pesos, datos, entornos, retroalimentación & ¿La capacidad proviene de parámetros estáticos o de un ciclo cerrado continuo? \\
Aplicaciones y despliegue & Usuarios, flujos de trabajo, distribución, ingresos & ¿Cómo obtiene el producto context de alta calidad y señales de verificación? \\
Seguridad y gobernanza & Normas, auditoría, responsabilidad, instituciones & ¿Quién define los límites y quién asume el costo de los fallos? \\
\bottomrule
\end{tabular}
\caption{El stack del curso no es una lista de empresas, sino una tabla de dependencias de recursos y de asignación de responsabilidades.}
\end{table}

\subsection{Declaración de intereses del ponente y posición histórica}

Para valorar la solidez de la evidencia de un curso sobre la industria, esta sección vuelve a situar en el panorama la línea de tiempo de participación del ponente. Las inversiones tempranas, la participación en consejos de administración, las colaboraciones en productos o la fundación de empresas aportan detalles difíciles de obtener en materiales públicos, pero también pueden hacer que ciertos temas reciban más atención. Un buen método de estudio no consiste en descartar todas las opiniones porque exista un conflicto de intereses, sino en distinguir entre observación de primera mano, narrativa corporativa, inferencia macro y juicio de valor, y buscar evidencia verificable para cada una.

\begin{figure}[H]
\centering
\includegraphics[width=0.94\textwidth]{images/00-16-45-background-disclosures.jpg}
\caption{Declaración del ponente, en forma de línea de tiempo, de sus inversiones y de los proyectos que fundó o en los que participó.}
\figsource{Video oficial de Stanford Online 00:16:45, `images/00-16-45-background-disclosures.jpg`.}
\end{figure}

\begin{warningbox}{Lectura de la figura: la declaración no es un descargo de responsabilidad decorativo}
La línea de tiempo indica al lector qué juicios pueden provenir de experiencia interna y cuáles solo pueden ser analogías externas. Por ejemplo, alguien que lleva mucho tiempo involucrado en empresas de modelos y en el mercado de capacidad de cómputo percibe con más facilidad la capital allocation y el deployment feedback, pero no necesariamente representa a todos los equipos de investigación, las instituciones públicas o los desarrolladores pequeños y medianos. Cuando más adelante se citen las opiniones del ponente sobre las context wars, los precios de las GPU y los arreglos institucionales, deben marcarse como juicios basados en la experiencia con una fuente identificada, y no como leyes objetivas sin sujeto.
\end{warningbox}

\subsection{Cuatro tipos de restricciones de capacidad}

Con el mapa de todo el stack, el curso agrupa los límites del progreso de la IA en cuatro grupos: context, compute, capital y culture. Este orden merece atención. El debate público suele tratar el compute como el único bien escaso, pero el cómputo solo se convierte en capacidad cuando entra en contacto con las tareas, los entornos y la retroalimentación adecuados; el capital determina la velocidad de expansión y la asunción de riesgos, y la cultura influye en qué está dispuesto a probar un equipo, qué fallos reconoce y qué valores incorpora a los productos y a las instituciones.

\begin{figure}[H]
\centering
\includegraphics[width=0.9\textwidth]{images/00-17-45-capability-bottlenecks.jpg}
\caption{Los cuatro tipos de restricciones al progreso de las capacidades de la IA: context, compute, capital y culture.}
\figsource{Video oficial de Stanford Online 00:17:45, `images/00-17-45-capability-bottlenecks.jpg`.}
\end{figure}

\begin{importantbox}{Cómo interactúan los cuatro tipos de restricciones}
\term{Context} determina qué tareas, herramientas y retroalimentación ve el sistema; \term{compute} aporta el presupuesto de entrenamiento e inferencia; \term{capital} convierte las expectativas de ingresos futuros en inversión en salas de servidores, chips y talento; \term{culture} determina cómo la organización comparte los errores, gestiona la seguridad, elige objetivos y coordina la colaboración. Cuando una capa es suficiente de momento, la presión aparece en las demás. Al diagnosticar un sistema hay que buscar la restricción marginal actual; no se puede reutilizar siempre la misma respuesta.
\end{importantbox}

\begin{knowledgebox}{Términos clave: bottleneck no es una etiqueta permanente}
Aquí, bottleneck designa el eslabón que más limita la producción total cuando se agrega una unidad de recurso en el momento actual. Cambia con la escala, la tarea, la región y el tiempo. Por ejemplo, a un equipo pequeño puede faltarle primero capital; a un gran laboratorio, primero electricidad e interconexión, y a un coding agent maduro, sobre todo entornos ejecutables y retroalimentación de usuarios reales. Si el curso cruza capas con frecuencia más adelante, es precisamente porque las restricciones se desplazan y también se amplifican entre sí.
\end{knowledgebox}

\subsection{Resumen de la sección}

La Great Transition significa que todo el stack técnico-económico está cambiando a la vez. El mapa de todo el stack aporta las relaciones de dependencia, la línea de tiempo de la declaración de intereses ayuda a juzgar la procedencia de las opiniones y los cuatro tipos de restricciones ofrecen un marco de diagnóstico dinámico. A continuación nos centramos en la primera restricción: cómo el context forma, mediante entornos, retroalimentación y aprendizaje por refuerzo, un ciclo cerrado de capacidades sostenible.
\section{RL y Context: por qué la capacidad surge del ciclo cerrado}

Si un modelo solo aprende de un corpus estático, puede absorber los patrones existentes, pero difícilmente sabrá de forma directa si una acción concreta tuvo éxito en una tarea real. El aprendizaje por refuerzo coloca al modelo en un entorno interactivo, donde intenta acciones, observa los resultados y actualiza su política según la reward. Aquí, \term{RL} es Reinforcement Learning, que en español suele traducirse como aprendizaje por refuerzo; \term{environment} es el mundo de la tarea en el que el modelo puede actuar y recibir retroalimentación sobre el estado o el resultado. Ambos amplían «más tokens» a «más experiencia verificada».

\begin{figure}[H]
\centering
\includegraphics[width=0.9\textwidth]{images/00-19-45-how-rl-works.jpg}
\caption{Ciclo cerrado del aprendizaje por refuerzo: el entorno entrega un estado, la política ejecuta una acción y la retroalimentación impulsa la actualización de la política.}
\figsource{Video oficial de Stanford Online 00:19:45, `images/00-19-45-how-rl-works.jpg`.}
\end{figure}

\begin{knowledgebox}{Lectura de la figura: separar los tres tipos de objetos}
En la figura, el agent o policy es el modelo que toma decisiones, el environment es el sistema externo que responde a las acciones y la reward es una retroalimentación comprimida sobre el resultado. Ninguno puede sustituir a los otros: un modelo más grande no equivale a un mejor entorno, un entorno más rico no garantiza que la reward sea correcta, y que la reward pueda calcularse no significa que represente el resultado que las personas realmente quieren. La capacidad del sistema surge de la calidad del ciclo cerrado, no de la escala aislada de alguno de sus componentes.
\end{knowledgebox}

Un objetivo de RL simplificado puede escribirse como
\[
J(\theta)=\mathbb{E}_{\tau\sim\pi_\theta}\left[\sum_{t=0}^{T}\gamma^t r_t\right].
\]
Aquí, $\pi_\theta$ denota la política con parámetros $\theta$, $\tau$ es una trajectory que surge de la interacción entre la política y el entorno, $r_t$ es la reward del paso $t$, $\gamma$ es el factor de descuento de los rendimientos futuros y $T$ es el instante en que termina la tarea. La fórmula no responde cómo debe diseñarse la reward ni garantiza que el entrenamiento sea estable; solo expresa como objeto de optimización «hacer más probables las trayectorias de alto rendimiento».

\begin{warningbox}{Nota de clase: el scaling empírico no es scaling ilimitado}
El ponente describe una y otra vez las scaling laws como empirical observations. Que el desempeño haya mejorado de forma sostenida al aumentar la escala de entrenamiento solo indica que existe una regularidad dentro del intervalo observado. El RL también puede estancarse por la calidad del entorno, fallas en el verificador, exploración insuficiente, cambios de distribución u objetivos mal alineados. La planeación de ingeniería puede extrapolar curvas, pero debe preparar al mismo tiempo las condiciones de falla y experimentos alternativos.
\end{warningbox}

\subsection{Basic AI Scaling Recipe}

Esta sección reduce el proceso de producción de las capacidades de los modelos modernos a cuatro nodos: el compute sostiene la research, la research produce models, los models llegan a los usuarios y las tareas mediante la inference, y el uso genera a su vez context y retroalimentación. Lo importante de este diagrama no es el número de flechas, sino la dirección del ciclo cerrado. El entrenamiento no encapsula el conocimiento en los pesos de una sola vez; el entorno después del despliegue, las llamadas a herramientas y los registros de éxitos y fracasos determinan lo que la siguiente ronda de investigación y entrenamiento podrá ver.

\begin{figure}[H]
\centering
\includegraphics[width=0.76\textwidth]{images/00-22-34-basic-ai-scaling-recipe.jpg}
\caption{Basic AI Scaling Recipe: el ciclo de compute, research, models, inference y context feedback.}
\figsource{Video oficial de Stanford Online 00:22:34, `images/00-22-34-basic-ai-scaling-recipe.jpg`.}
\end{figure}

\begin{importantbox}{Lectura de la figura: la flecha del ciclo que más se pasa por alto}
La ruta de compute a model es fácil de cuantificar: el número de GPU, los FLOPs de entrenamiento y el tamaño del modelo tienen cifras. La ruta de regreso de inference a context feedback es más difícil: exige saber en qué flujo de trabajo usan los usuarios el modelo, cuándo falla, si la falla puede verificarse, si hay autorización para que los datos regresen y si la retroalimentación representa a la población objetivo. El ponente coloca el context en primer lugar porque esta flecha suele determinar si el compute posterior tendrá usos de alto valor.
\end{importantbox}

\subsection{El context no es solo un prompt más largo}

Aquí, el alcance del context es mucho mayor que la ventana de conversación. Incluye la descripción de la tarea, los documentos privados, el código fuente, los permisos de herramientas, el entorno de ejecución, las decisiones históricas, el comportamiento de los usuarios y los resultados verificables. Aun con los mismos pesos, un modelo que difiere en si puede acceder a la terminal, a las pruebas, a los registros de despliegue y al conocimiento de la organización se comporta como un sistema completamente distinto. Por eso Frontier Systems considera el context un activo de infraestructura: hay que recopilarlo, autorizarlo, depurarlo, conectarlo y mantenerlo de forma continua.

\begin{figure}[H]
\centering
\includegraphics[width=0.9\textwidth]{images/00-24-04-context-environments-flywheel.jpg}
\caption{Context/Environments flywheel: integrarse al flujo de trabajo, observar las fallas, mejorar el entorno y el modelo, y ampliar el uso.}
\figsource{Video oficial de Stanford Online 00:24:04, `images/00-24-04-context-environments-flywheel.jpg`.}
\end{figure}

\begin{knowledgebox}{Lectura de la figura: tres preguntas determinan si el flywheel funciona}
Primero, si el sistema entra en contacto con tareas reales de alto valor o solo gira sobre datos de demostración; segundo, si el éxito y el fracaso pueden verificarse de forma confiable o solo se cuenta con una satisfacción difusa; tercero, si el regreso de los datos cuenta con la autorización del usuario y respeta los límites de seguridad. El flywheel funciona solo cuando el crecimiento del uso, la calidad de la retroalimentación y la mejora del modelo se refuerzan entre sí. Aumentar únicamente el volumen de llamadas puede amplificar el ruido, los riesgos de privacidad y la automatización de errores.
\end{knowledgebox}

\begin{table}[H]
\centering
\begin{tabular}{L{0.18\textwidth}L{0.32\textwidth}L{0.4\textwidth}}
\toprule
Término & Problema que resuelve & Significado en esta clase \\
\midrule
Prompt context & Qué información necesita la solicitud actual & Conversación, documentos y restricciones inmediatas \\
Tool context & Qué acciones puede realizar el modelo & API, terminal, búsqueda, bases de datos y permisos \\
Environment & Dónde se ejecuta la acción y produce resultados & Repositorio de código, navegador, laboratorio o sistema de negocio \\
Feedback & Cómo juzgar el efecto de una acción & Pruebas, modificaciones del usuario, ingresos, latencia o incidentes de seguridad \\
Flywheel & Por qué las mejoras pueden acumularse de forma continua & El uso genera retroalimentación de alta calidad, y esta mejora el modelo y el producto \\
\bottomrule
\end{tabular}
\caption{Glosario: descomponer el context genérico en componentes que pueden diseñarse y auditarse.}
\end{table}

\subsection{Resumen de la sección}

El RL ofrece un ciclo cerrado de aprendizaje política–entorno–retroalimentación, y la Basic Scaling Recipe sitúa ese ciclo dentro del proceso full-stack de compute, research, model e inference. Lo decisivo del context no es su longitud, sino la calidad de las tareas, las herramientas, los entornos y las señales de verificación. Y precisamente porque el context de alto valor es escaso, las empresas de modelos y las empresas de aplicaciones compiten intensamente por los puntos de acceso, los datos y los flujos de trabajo.
\section{Context Loop Wars: puntos de entrada, herramientas y retorno de datos}

Cuando un coding agent, un agent de oficina o un agent sectorial entra en un flujo de trabajo real, la empresa que controla el punto de entrada del usuario puede observar la distribución de las tareas, la empresa dueña del modelo puede controlar la capacidad y el precio, y la empresa dueña de las herramientas y del entorno define los límites de acción. Estas tres posiciones no siempre las ocupa la misma empresa, por lo que las relaciones de colaboración se convierten rápidamente en competencia. Esta sección trata los casos de Windsurf, OpenAI, Anthropic y Mistral que eligió el expositor como un análisis de estructura de sistemas, no como chisme empresarial.

\subsection{Qué reveló una disputa por una adquisición}

El expositor usa los cambios del mercado de coding agents en 2025 para mostrar el valor del context loop. Windsurf se ubica en el punto de entrada del flujo de trabajo de los desarrolladores, los proveedores de modelos aportan la capacidad de base, y el código de los usuarios junto con la retroalimentación de la ejecución forman un entorno de alto valor. En cuanto una de las partes teme que la otra interiorice los datos de interacción, la distribución o la relación con el producto, la colaboración original vía API da paso a restricciones de acceso, negociaciones de adquisición y despliegue de modelos alternativos.

\begin{figure}[H]
\centering
\includegraphics[width=0.94\textwidth]{images/00-27-34-context-loop-wars.jpg}
\caption{Primera etapa de las Context Loop Wars: el conflicto en torno al punto de entrada del flujo de trabajo de Windsurf y al acceso a los modelos.}
\figsource{Video oficial de Stanford Online 00:27:34, `images/00-27-34-context-loop-wars.jpg`.}
\end{figure}

\begin{knowledgebox}{Lectura de la figura: no hay que fijarse solo en el precio de la adquisición}
Lo más importante son cuatro flujos: por dónde entran las solicitudes del usuario, dónde se ejecutan el código y las herramientas, qué empresa aporta el modelo que hace la inferencia y a qué sistema de datos regresan los resultados y las modificaciones. La empresa del punto de entrada controla la demanda y la distribución; la empresa del modelo controla la capacidad y la oferta; la empresa de nube o de herramientas controla la infraestructura de ejecución. Una transacción solo reorganiza el control de estos cuatro flujos. Una vez entendida la estructura, el método de análisis sigue siendo válido aunque cambien las empresas concretas.
\end{knowledgebox}

\begin{figure}[H]
\centering
\includegraphics[width=0.94\textwidth]{images/00-29-04-context-loop-wars-cont.jpg}
\caption{Continuación de las Context Loop Wars: modelos, nube, herramientas de desarrollo y talento se recombinan dentro del mismo ciclo cerrado.}
\figsource{Video oficial de Stanford Online 00:29:04, `images/00-29-04-context-loop-wars-cont.jpg`.}
\end{figure}

La segunda figura amplía una transacción aislada a una relación de ecosistema. Al leerla conviene comparar qué le falta a cada participante: a los laboratorios de modelos pueden faltarles una distribución estable y entornos de alta calidad; a las herramientas de desarrollo, modelos propios y control del costo de inferencia a largo plazo; las plataformas de nube pueden querer fijar la demanda de cómputo dentro de su propia infraestructura, y el talento y los equipos de investigación se desplazan según dónde estén el capital y los datos. La llamada context war es, en esencia, una disputa por un ciclo cerrado de mejora sostenible, no por un lote de registros de chat estáticos.

\begin{warningbox}{Nota de clase: los casos ilustran mecanismos, no garantizan que la cronología siga vigente}
El estado de las fusiones y adquisiciones, las políticas de API y las capacidades de los productos cambian. Estas notas conservan los casos históricos que el expositor usó en esa sesión para explicar el valor del context. El lector debe verificar los hechos más recientes y después usar el marco «punto de entrada—entorno—modelo—retroalimentación—distribución» para juzgar la nueva situación; las conexiones entre empresas de las diapositivas no deben tomarse como un mapa permanente de la industria.
\end{warningbox}

\subsection{Mistral: el context clave también puede ser un entorno soberano}

El context no pertenece solo a un producto SaaS. El idioma, el derecho, las compras del sector público, los requisitos de seguridad nacional, la residencia de los datos y los ecosistemas industriales locales también forman entornos difíciles de replicar de una región a otra. El expositor explica la motivación detrás de la creación de Mistral a partir de las necesidades europeas de soberanía y de pesos abiertos: cuando las capacidades clave de los modelos dependen por completo de proveedores externos, a las instituciones locales les resulta difícil controlar los datos, los costos, la capacidad de auditoría y su poder de negociación a largo plazo.

\begin{figure}[H]
\centering
\includegraphics[width=0.94\textwidth]{images/00-30-00-mistral-critical-contexts.jpg}
\caption{El caso de Mistral: los idiomas, la regulación, las instituciones públicas y las necesidades industriales de Europa constituyen critical contexts/environments.}
\figsource{Video oficial de Stanford Online 00:30:00, `images/00-30-00-mistral-critical-contexts.jpg`.}
\end{figure}

\begin{knowledgebox}{Lectura de la figura: el significado sistémico de la sovereign AI}
La sovereign AI no consiste solo en «entrenar un modelo nacional». Abarca, como mínimo, datos de idioma y cultura locales, licencias de modelo aceptables, residencia de los datos, control de la cadena de suministro, interfaces de despliegue para el sector público, talento y cómputo, así como la capacidad de mantener el servicio en una crisis. Los pesos abiertos pueden mejorar la posibilidad de inspección y las opciones de despliegue, pero no resuelven por sí solos el costo de entrenamiento, la responsabilidad en seguridad ni la fragmentación del ecosistema. Al final, las aspiraciones de soberanía tienen que concretarse en infraestructura operable y en arreglos institucionales.
\end{knowledgebox}

\begin{importantbox}{Las tres escaseces del context}
El context de producto es escaso porque los flujos de trabajo reales y la confianza de los usuarios son difíciles de replicar; el context organizacional es escaso porque el conocimiento privado, los permisos y las decisiones históricas no pueden recopilarse públicamente; el context soberano es escaso porque el idioma, el derecho y el interés público son territoriales. Cuanto más se parezcan las capacidades de los modelos, más probable será que quien pueda entrar de forma segura en estos entornos obtenga la siguiente ronda de retroalimentación de alto valor.
\end{importantbox}

\subsection{Resumen de la sección}

Las context loop wars concretan la abstracta «ventaja de datos» como el control sobre el punto de entrada, las herramientas, el entorno de ejecución, la oferta de modelos y el retorno de la retroalimentación. El caso de Windsurf muestra el ciclo cerrado de producto; el caso de Mistral, los entornos soberanos y públicos. La siguiente sección plantea otra pregunta: aun con el entorno y la retroalimentación, ¿puede el RL seguir escalando de forma sostenida, o se topará con la frontera entre lo medible y lo no medible?
\section{Frontier RL Flywheel y límites de capacidad}

Con un entorno real, el sistema puede generar tareas, intentar acciones, verificar resultados y luego usar las trayectorias exitosas para actualizar el modelo. El ponente llama a este proceso recursive self-improvement, pero el término se malinterpreta con facilidad, como si el modelo se separara por completo de las personas y se mejorara de manera autónoma e ilimitada. Una explicación más precisa del sistema es esta: el modelo ayuda a ampliar la experimentación y la producción de datos, mientras que las personas y la infraestructura siguen encargándose de elegir los entornos, definir la verificación, controlar el despliegue, atender las anomalías y decidir qué mejoras merecen pasar a la siguiente ronda.

\begin{figure}[H]
\centering
\includegraphics[width=0.94\textwidth]{images/00-35-15-frontier-rl-flywheel.jpg}
\caption{Frontier RL flywheel: parte de hipótesis y entornos, pasa por el intento, la verificación y el aprendizaje, y genera un sistema más fuerte para la siguiente ronda.}
\figsource{Video oficial de Stanford Online, 00:35:15, `images/00-35-15-frontier-rl-flywheel.jpg`.}
\end{figure}

\begin{knowledgebox}{Lectura de la figura: la mejora recursiva necesita anclas externas}
A la izquierda está la generación de tareas, datos o hipótesis científicas; al centro, el modelo actúa en el entorno y obtiene resultados verificables; a la derecha, la experiencia útil se incorpora de vuelta a la política, las herramientas o la base de conocimiento. Las anclas externas del ciclo incluyen pruebas unitarias, experimentos físicos, aceptación por parte del usuario, reglas de seguridad y límites de costo. Sin anclas confiables, el sistema podría limitarse a optimizar métricas sustitutas y caer en un autorreforzamiento que parece rápido, pero que en realidad se aleja del objetivo.
\end{knowledgebox}

\subsection{Límite uno: cuánto más puede escalar en la práctica}

El límite empírico se ocupa de curvas observables: al aumentar los rollouts, el número de entornos, la dificultad de las tareas y la capacidad de cómputo para entrenamiento, ¿sigue subiendo la tasa de éxito?, ¿cuándo empiezan a disminuir los rendimientos marginales?, ¿aparece colapso de modos o reward hacking durante el entrenamiento? Esta pregunta puede responderse poco a poco mediante experimentos. La respuesta también varía según el dominio: las pruebas de código, las respuestas matemáticas, los experimentos de materiales y la escritura abierta tienen costos de verificación completamente distintos.

\begin{figure}[H]
\centering
\includegraphics[width=0.9\textwidth]{images/00-37-15-limits-to-rl.jpg}
\caption{El ponente plantea dos tipos de límites de RL: la meseta empírica y el límite filosófico relativo a los objetivos y los valores humanos.}
\figsource{Video oficial de Stanford Online, 00:37:15, `images/00-37-15-limits-to-rl.jpg`.}
\end{figure}

\begin{warningbox}{Lectura de la figura: los dos tipos de pregunta no se responden con la misma curva}
«¿Sigue subiendo la exactitud con más capacidad de cómputo?» es una pregunta empírica, para la que pueden diseñarse experimentos controlados; «¿qué se debe optimizar?, ¿quién tiene derecho a definir el éxito?, ¿qué capacidades no deben desplegarse?» son preguntas normativas y de gobernanza. Estas últimas no desaparecen por sí solas porque el benchmark siga subiendo. Mezclar ambas hace que la curva técnica sustituya a la elección de valores, y también que el debate sobre valores ignore la evidencia de ingeniería real.
\end{warningbox}

\subsection{Límite dos: de quién es el objetivo que representa la reward}

El límite filosófico termina por convertirse en una interfaz del sistema. Quién elige las tareas de entrenamiento, quién escribe el verifier, quién puede modificar la reward, quién asume los fallos y si el usuario sabe qué está optimizando el sistema: ninguna de estas es una cuestión accesoria que se atienda después del entrenamiento. La ventaja de RL es que incorpora la retroalimentación al aprendizaje; su riesgo proviene de ese mismo punto: si el criterio de retroalimentación es demasiado estrecho, manipulable o poco representativo, la optimización amplificará el sesgo hasta convertirlo en una conducta estable.

\begin{table}[H]
\centering
\begin{tabular}{L{0.2\textwidth}L{0.31\textwidth}L{0.39\textwidth}}
\toprule
Tipo de límite & Señal observable & Forma de respuesta \\
\midrule
Límite de cómputo & Rendimiento de rollouts, duración del entrenamiento, energía y costo & Optimización del sistema, eficiencia de muestreo, presupuesto de recursos \\
Límite de entorno & Cobertura insuficiente de tareas, desfase entre simulación y realidad & Ampliar los entornos reales, verificación por capas, auditoría humana \\
Límite de verificación & Reward hacking, filtración de pruebas, distorsión de métricas sustitutas & Varios verifiers, pruebas ocultas, repositorio de casos de fallo \\
Límite organizacional & Responsabilidad difusa, presión de lanzamiento que se impone a las señales de seguridad & Autoridad de decisión clara, análisis posterior de incidentes, condiciones de detención \\
Límite de valores & Conflicto entre grupos sobre la definición de éxito & Diseño participativo, restricciones institucionales, mecanismos de apelación \\
\bottomrule
\end{tabular}
\caption{Descomposición de «los límites de RL» en problemas distintos que pueden medirse y gobernarse.}
\end{table}

\begin{knowledgebox}{Nota de clase: primero, recorrer el ciclo por cuenta propia}
El ponente anima a los estudiantes a entender RL mediante el problem set, en lugar de juzgarlo solo a partir de narrativas generales. Definir uno mismo el entorno, escribir la reward y observar cómo la política aprovecha los huecos pone en evidencia muy pronto la distancia entre «un objetivo que parece claro» y «un objetivo que puede calcularse de forma confiable». Esta práctica también es preparedness: saber cuándo conviene confiar en la curva y cuándo hay que revisar los datos y el verifier.
\end{knowledgebox}

\subsection{Resumen de la sección}

El Frontier RL flywheel puede ampliar la producción de experiencia verificable, pero cada ronda depende de los entornos, el verifier, las decisiones humanas y la infraestructura. El límite empírico puede explorarse con curvas y experimentos; los límites filosóficos y de gobernanza requieren responsabilidades claras y deliberación pública. A continuación, la perspectiva pasa del ciclo de entrenamiento individual al mercado en su conjunto, para observar si el gasto de capital, la adopción de software y el precio del cómputo respaldan la narrativa de la «gran transformación».
\section{Scaling macro: de las curvas de los modelos a las señales de capital y de precios}

La expansión de la IA suele resumirse en una sola curva, pero el sistema industrial produce al menos cuatro tipos de señales a la vez: la señal técnica, que muestra cómo cambia la capacidad de los modelos según los recursos de entrenamiento; la señal de producto, que es el crecimiento en el uso de las herramientas; la señal de capital, que es la inversión en centros de datos y chips; y la señal de oferta y demanda, que dan los precios de renta de GPU. Estas señales pueden respaldarse entre sí, pero no son lo mismo. El progreso técnico no garantiza ingresos comerciales, el crecimiento de los ingresos no garantiza que el capital se asigne con eficacia y la inversión de capital tampoco garantiza que cada pieza de hardware se aproveche de manera eficiente.

\subsection{La previsibilidad de la capacidad de los modelos}

El valor de ingeniería de la scaling law está en la gestión del presupuesto y del riesgo: con experimentos anteriores, un equipo puede estimar cómo cambiará la pérdida si aumentan los datos, los parámetros o el cómputo, y con eso decidir el tamaño del siguiente entrenamiento. Mientras más lejos llega el intervalo de predicción, mayor es el error que introducen los cambios de arquitectura, la calidad de los datos, el posentrenamiento y la deriva de la evaluación. Las gráficas macro sirven para mostrar la dirección a largo plazo, pero no reemplazan la auditoría de cada modelo, cada tarea y cada recipe de entrenamiento.

\begin{figure}[H]
\centering
\includegraphics[width=0.94\textwidth]{images/00-44-00-macro-ai-scaling.jpg}
\caption{El ponente usa una gráfica macro de la expansión de los modelos para mostrar que el AI scaling presenta regularidades observables en varios periodos históricos.}
\figsource{Video oficial de Stanford Online 00:44:00, `images/00-44-00-macro-ai-scaling.jpg`.}
\end{figure}

\begin{knowledgebox}{Lectura de la figura: distinguir puntos observados, línea de ajuste y supuestos sobre el futuro}
Los puntos observados provienen de proyectos de entrenamiento ya publicados o estimados; la línea de ajuste resume la tendencia pasada, y el tramo que se extiende hacia la derecha incluye supuestos sobre el futuro. Conviene revisar primero las unidades de los ejes y la fuente de los datos, y después preguntar si las generaciones de modelos son comparables, si el posentrenamiento se tomó en cuenta y si la evaluación cambió. Un buen ajuste solo demuestra que los datos elegidos muestran una regularidad clara en ese sistema de coordenadas; no demuestra que la siguiente generación de sistemas vaya a crecer con la misma pendiente.
\end{knowledgebox}

\subsection{Adopción del producto: qué indica el crecimiento de Claude Code}

Las señales de producto están más cerca de la demanda real que los benchmark. El ponente muestra el crecimiento de los commits diarios de Claude Code a lo largo de doce meses para ilustrar que los coding agent se están incorporando rápidamente al proceso de desarrollo. Este indicador significa, como mínimo, que los usuarios están dispuestos a subir al repositorio código generado o asistido por el modelo, pero por sí solo no responde sobre la calidad del código, el costo del retrabajo, los incidentes de seguridad, la retención de clientes de pago ni la productividad neta de distintos equipos.

\begin{figure}[H]
\centering
\includegraphics[width=0.94\textwidth]{images/00-46-02-claude-code-growth.jpg}
\caption{Los commits diarios de Claude Code crecieron rápidamente en doce meses, lo que muestra la velocidad con la que los productos de agent entran en flujos de trabajo reales.}
\figsource{Video oficial de Stanford Online 00:46:02, `images/00-46-02-claude-code-growth.jpg`.}
\end{figure}

\begin{importantbox}{Lectura de la figura: el volumen de adopción no es una conclusión sobre productividad}
El número de commits puede demostrar que la intensidad de uso aumenta, pero no demuestra directamente que cada commit eleve la producción neta. Una evaluación completa también requiere merge rate, review time, tasa de defectos, tasa de reversiones, cobertura de desarrolladores y costo de mantenimiento a largo plazo. El crecimiento del producto sigue siendo importante porque amplía las fuentes de context y de retroalimentación; solo que esa retroalimentación debe pasar por filtros de calidad y de permisos antes de alimentar la siguiente ronda de mejoras del modelo o de las herramientas.
\end{importantbox}

\subsection{El gasto de capital de la Great Transition}

Cuando la demanda de aplicaciones, el tamaño de los modelos y las expectativas de capacidad de cómputo suben al mismo tiempo, los proveedores de nube y las empresas de infraestructura compran por adelantado terrenos, energía eléctrica, redes y aceleradores. \term{Capex} es capital expenditure, es decir, el gasto de capital destinado a activos de largo plazo; a diferencia de los gastos operativos cotidianos, suele tener ciclos de construcción largos, recuperación lenta y una salida difícil. Si los pronósticos de demanda de IA se desvían, tanto la sobreconstrucción como la escasez de oferta pueden prolongarse durante años.

\begin{figure}[H]
\centering
\includegraphics[width=0.94\textwidth]{images/00-46-45-great-transition.jpg}
\caption{El ponente describe el aumento reciente del gasto de capital en infraestructura como una señal macro de la Great Transition.}
\figsource{Video oficial de Stanford Online 00:46:45, `images/00-46-45-great-transition.jpg`.}
\end{figure}

\begin{knowledgebox}{Lectura de la figura: el capex es una apuesta, no un valor ya realizado}
La curva de gasto de capital indica que las empresas creen que habrá suficientes ingresos de software y valor estratégico para sostener la inversión en hardware. Una vez construido el centro de datos, su valor depende además de la tasa de utilización, la eficiencia de los modelos, el precio de la electricidad, la capacidad de interconexión y la demanda de los clientes. Es razonable tomar el capex como una expresión cuantitativa de la confianza en la demanda; tomarlo como prueba de que la productividad o el beneficio social ya se materializaron implica saltarse una cadena que todavía no se ha verificado.
\end{knowledgebox}

\subsection{Precios de renta de GPU: el termómetro de un mercado heterogéneo de cómputo}

El precio resulta atractivo porque condensa en un solo número la oferta, la demanda, la región, la energía eléctrica, la generación del hardware y el plazo del contrato. El ponente muestra la evolución de los precios de renta de GPU observada por AMP para ilustrar que la escasez no aumenta de forma monótona: la entrada de nueva oferta, la depreciación de las tarjetas antiguas, las mejoras en la eficiencia del software y los cambios en la estructura de la demanda modifican el precio. Más importante aún, las distintas GPU no son una misma mercancía; el precio de referencia debe leerse junto con la memoria de video, la interconexión, la confiabilidad, el tiempo de disponibilidad y el tamaño del clúster.

\begin{figure}[H]
\centering
\includegraphics[width=0.94\textwidth]{images/00-47-45-gpu-rental-prices.jpg}
\caption{Los precios de renta de GPU cambian según la generación del hardware y la etapa del mercado, y reflejan el efecto conjunto de la oferta y la demanda, la depreciación y la sustituibilidad.}
\figsource{Video oficial de Stanford Online 00:47:45, `images/00-47-45-gpu-rental-prices.jpg`.}
\end{figure}

\begin{warningbox}{Lectura de la figura: el precio por GPU-hour sigue sin ser el precio de una mercancía uniforme}
Aunque todo se exprese por hora, una sola tarjeta, un nodo de ocho tarjetas y un clúster grande sin bloqueo tienen valores completamente distintos; la renta de corto plazo, la reservada, la interrumpible y los contratos de largo plazo tampoco se pueden comparar directamente. El entrenamiento depende además de la topología de red, el rendimiento del almacenamiento, la pila de software y la recuperación ante fallas. La gráfica de precios puede mostrar una tendencia, pero sin las especificaciones y las condiciones del servicio no puede demostrar que «el cómputo ya es barato» ni que «cierta tarjeta conviene más con seguridad».
\end{warningbox}

\subsection{Resumen de la sección}

Las curvas de los modelos, la adopción de productos, el gasto de capital y los precios de renta responden, respectivamente, a preguntas de tecnología, demanda, inversión y oferta y demanda. Ponerlas lado a lado permite observar si la transición ocurre a través de varias capas, pero no permite fundir los cuatro tipos de señales en una cadena causal necesaria. La siguiente sección recurre a los ciclos históricos de infraestructura para buscar patrones reutilizables y, al mismo tiempo, deja claros los límites de la analogía.
\section{Ciclos históricos de infraestructura: escasez, sobreconstrucción y estandarización}

El ponente sitúa el compute actual dentro de una historia más larga de la infraestructura: el acero, el ferrocarril, la electricidad, la fibra óptica, los materiales semiconductores y las tecnologías energéticas pasaron por un auge de la demanda, una afluencia de capital, una expansión de la oferta, una caída de precios, la consolidación de empresas y la formación de estándares. La analogía histórica sirve para generar preguntas, como «dónde habrá sobreconstrucción», «qué interfaces se estandarizarán» o «quién asume los costos hundidos»; no permite predecir calendarios directamente, porque cada recurso difiere en sus propiedades físicas, su estructura regulatoria y su ritmo de innovación.

\subsection{Gilded Age: la infraestructura y el poder se expanden juntos}

La Gilded Age de finales del siglo XIX suele usarse para describir una etapa de rápida expansión del ferrocarril, el acero y las finanzas, acompañada de concentración del poder de mercado, conflictos laborales y una regulación pública cada vez más intensa. El ponente la utiliza para recordar a los estudiantes que el auge de una tecnología de propósito general no solo trae más capacidad productiva, sino que también redistribuye el poder de negociación. Por ello, la cuestión central de la infraestructura de IA no es solo «cuánto puede construirse», sino también quién controla el acceso, quién fija los precios y quién puede seguir innovando en tiempos de escasez.

\begin{figure}[H]
\centering
\includegraphics[width=0.88\textwidth]{images/00-51-00-gilded-age.jpg}
\caption{La Gilded Age sirve como punto de entrada histórico de esta clase: el auge de la infraestructura y la concentración del poder suelen aparecer al mismo tiempo.}
\figsource{Video oficial de Stanford Online 00:51:00, `images/00-51-00-gilded-age.jpg`.}
\end{figure}

\begin{knowledgebox}{Lectura de la figura: en una analogía histórica, primero se comparan los mecanismos}
Los mecanismos comparables incluyen los altos costos fijos, los efectos de red, el control de los recursos, la integración vertical, la discriminación de precios y la presión de la regulación pública. Las partes que no admiten una analogía directa incluyen la velocidad de actualización tecnológica, las cadenas de suministro globales, el costo de replicar software y los mercados de capital modernos. Tomar «hoy es una nueva Gilded Age» como conclusión es demasiado burdo; tomarlo como un recordatorio para examinar la concentración del mercado y el interés público tiene más valor didáctico.
\end{knowledgebox}

\subsection{El acero: cómo las mejoras de proceso cambian el precio de un recurso}

Innovaciones de proceso como el Bessemer process redujeron de forma notable el costo de producción del acero, lo que permitió que el ferrocarril, la construcción y la maquinaria se expandieran a mayor escala. Este caso corresponde, en la IA, a la eficiencia de los modelos, el diseño de chips y la optimización de sistemas: el crecimiento de la demanda no tiene que satisfacerse únicamente aumentando la oferta de materias primas, porque el propio proceso de producción también modifica el costo unitario. A su vez, la caída de precios libera nueva demanda, de modo que el consumo total sigue aumentando.

\begin{figure}[H]
\centering
\includegraphics[width=0.94\textwidth]{images/00-52-00-bessemer-steel-prices.jpg}
\caption{Tras el proceso Bessemer, el precio del acero bajó durante un largo periodo: la eficiencia productiva cambió la accesibilidad de la infraestructura.}
\figsource{Video oficial de Stanford Online 00:52:00, `images/00-52-00-bessemer-steel-prices.jpg`.}
\end{figure}

\begin{importantbox}{Lectura de la figura: la caída del costo unitario puede ampliar la demanda total}
Cuando el acero se abarató, la sociedad no dejó de usarlo, sino que lo llevó a más vías férreas, puentes y edificios. La eficiencia de cómputo también puede producir un rebound effect similar: si cada inferencia es más barata, surgen más agent, tareas más largas e interacciones más frecuentes. Al estimar la demanda total de energía y de capacidad de cómputo, no basta con ver la reducción del costo por tarea; también hay que estimar la elasticidad de la demanda que generan los nuevos usos.
\end{importantbox}

\subsection{La fibra óptica: la sobreconstrucción también puede dejar capacidad pública}

La gran cantidad de fibra óptica tendida durante la burbuja de internet provocó a corto plazo quiebras y depreciación de activos, pero dejó capacidad de bajo costo para la banda ancha y los servicios en la nube posteriores. Esta historia suele usarse para justificar la expansión de los centros de datos de IA, pero entre ambos casos hay diferencias importantes: la fibra oscura puede permanecer inactiva durante mucho tiempo, mientras que las GPU se deprecian más rápido, y la renovación generacional de los chips, el encapsulado y la compatibilidad del software acortan su vida útil efectiva. Que la sobreconstrucción tenga valor público a largo plazo depende de que los activos puedan reutilizarse.

\begin{figure}[H]
\centering
\includegraphics[width=0.94\textwidth]{images/00-53-00-fiber-optic-overbuild.jpg}
\caption{El precio de la fibra óptica cayó rápidamente tras la sobreconstrucción, y la capacidad ociosa se convirtió después en la base de la expansión de internet.}
\figsource{Video oficial de Stanford Online 00:53:00, `images/00-53-00-fiber-optic-overbuild.jpg`.}
\end{figure}

\begin{warningbox}{Lectura de la figura: ``primero sobreconstruir y luego ver'' no es una estrategia incondicional}
La vida útil de los activos, los costos de mantenimiento, la compatibilidad con los estándares y la dificultad de redistribuirlos determinan si la sobreconstrucción puede convertirse en capacidad pública. Si un centro de datos está atado a una sola red eléctrica, a una interconexión cerrada o a chips que se vuelven obsoletos con rapidez, su valor residual puede ser muy bajo. La planeación de infraestructura debe considerar al mismo tiempo la modularity, la posibilidad de actualización y las interfaces abiertas, en lugar de comparar solo la escala máxima de construcción.
\end{warningbox}

\subsection{La DRAM: por qué la oferta de semiconductores es cíclica}

\term{DRAM} significa Dynamic Random-Access Memory; en español, memoria dinámica de acceso aleatorio. Almacena cada bit en un capacitor, requiere refrescarse periódicamente y se usa ampliamente como memoria principal en servidores y computadoras personales. La inversión en plantas de fabricación de obleas tiene ciclos largos y costos fijos altos, mientras que la demanda varía con la coyuntura de los dispositivos y los centros de datos, por lo que los precios suelen oscilar entre la escasez y el exceso. La memoria de video y el encapsulado avanzado que usan los aceleradores de IA están sujetos a restricciones similares de planeación de capacidad.

\begin{figure}[H]
\centering
\includegraphics[width=0.94\textwidth]{images/00-53-30-dram-cycles.jpg}
\caption{El ciclo de la industria de la DRAM muestra cómo se repiten el retraso en la ampliación de capacidad, el ajuste de inventarios y la volatilidad de precios.}
\figsource{Video oficial de Stanford Online 00:53:30, `images/00-53-30-dram-cycles.jpg`.}
\end{figure}

\begin{knowledgebox}{Lectura de la figura: cómo el retraso de la oferta genera ciclos}
Cuando la demanda aumenta, primero suben los precios y los fabricantes amplían su capacidad en consecuencia; la nueva capacidad entra de golpe tras la construcción y el arranque gradual, el mercado puede pasar al exceso de oferta, y la caída de precios obliga a las empresas a recortar la inversión. Cuando la demanda se recupera, los recortes de la ronda anterior vuelven a estrechar la oferta. Las GPU, la HBM y el encapsulado avanzado no son idénticos a la DRAM, pero en todos los casos hay que analizar en un mismo modelo los ciclos largos de construcción y los pronósticos de demanda de corto plazo.
\end{knowledgebox}

\subsection{El uranio y la energía: la política, los permisos y la geopolítica cambian la curva}

La expansión de la capacidad de cómputo termina dependiendo de la energía, por lo que esta sección sigue el rastro desde los chips y las salas de servidores hasta la generación eléctrica, el combustible y la red eléctrica. El mercado del uranio muestra que el precio de un recurso no depende solo de la oferta de las minas y de la demanda, sino también de la política nuclear, los ciclos de permisos, los inventarios, la geopolítica y la aceptación pública. Aunque una fuente de energía sea técnicamente adecuada para un suministro estable, que pueda llegar a tiempo a los centros de datos sigue dependiendo de la interconexión con la red, el financiamiento, la construcción y la coordinación regulatoria. Al leer la figura hay que distinguir entre el precio del mineral, la electricidad que puede ponerse en operación y la capacidad realmente disponible para un complejo específico; entre los tres hay años de retraso por construcción y aprobación.

\begin{figure}[H]
\centering
\includegraphics[width=0.94\textwidth]{images/00-54-15-uranium-cycle.jpg}
\caption{El ciclo del precio del uranio refleja cómo los recursos, la política, los inventarios y la demanda energética moldean en conjunto el costo de la infraestructura.}
\figsource{Video oficial de Stanford Online 00:54:15, `images/00-54-15-uranium-cycle.jpg`.}
\end{figure}

\begin{knowledgebox}{Lectura de la figura: el precio no es una medida pura de la escasez física}
La cantidad de yacimientos solo determina una parte de la oferta; los permisos de extracción, la capacidad de conversión y enriquecimiento, los contratos de largo plazo, las reservas estratégicas y los planes de construcción de centrales nucleares modifican la oferta disponible. En el debate sobre la energía para la IA también hay que evitar tratar las cifras en GW como electricidad disponible en cualquier momento: la ubicación, el horario, la estabilidad, la capacidad de transmisión y las restricciones de carbono determinan si puede abastecer a un clúster concreto.
\end{knowledgebox}

\subsection{Alternative assets: detrás de una misma forma de onda puede haber mecanismos distintos}

El ponente reúne los ciclos de precios de varias clases de alternative assets para buscar una estructura temporal común que va de la escasez al entusiasmo, a la ampliación de capacidad y a la caída. Esta comparación puede ayudar a los equipos de inversión y construcción a delimitar un rango de escenarios, pero también puede crear con facilidad la ilusión de que «todos los ciclos son iguales». Curvas parecidas pueden deberse a causas distintas, y su duración también depende de la vida útil de los activos, la regulación y los sustitutos.

\begin{figure}[H]
\centering
\includegraphics[width=0.94\textwidth]{images/00-55-45-alternatives-timeline.jpg}
\caption{Comparación de los ciclos de varios alternative assets, para discutir el discurso de la escasez, la afluencia de capital y la respuesta de la oferta.}
\figsource{Video oficial de Stanford Online 00:55:45, `images/00-55-45-alternatives-timeline.jpg`.}
\end{figure}

\begin{warningbox}{Lectura de la figura: una forma parecida no implica la misma causa}
Que un precio suba y luego baje puede deberse a la ampliación de capacidad, al colapso de la demanda, a cambios de política, a la sustitución tecnológica o al desapalancamiento financiero. Si solo se hace una comparación visual, se pasan por alto los plazos de entrega, las diferencias de calidad y los efectos de red de cada activo. Al usar una analogía histórica, primero hay que explicitar el mecanismo y después comprobar si el mercado actual de compute cumple las mismas condiciones.
\end{warningbox}

\begin{table}[H]
\centering
\begin{tabular}{L{0.14\textwidth}L{0.21\textwidth}L{0.23\textwidth}L{0.25\textwidth}}
\toprule
Caso & Mecanismo transferible & Implicación para el compute & Lo que no admite analogía directa \\
\midrule
Acero & La innovación de proceso reduce el costo unitario & La eficiencia del software y de los chips cambia la accesibilidad & La demanda digital crece más rápido \\
Fibra óptica & Tras la sobreconstrucción, los activos depreciados se reutilizan & La capacidad construida por adelantado puede sostener aplicaciones posteriores & Las GPU se deprecian y cambian de generación más rápido \\
DRAM & Desfase entre ciclos largos de ampliación y ciclos cortos de demanda & Los chips, el encapsulado y los inventarios fluctúan & Productos más heterogéneos y cadena de suministro más compleja \\
Uranio & La política, los permisos y la geopolítica moldean la oferta & La electricidad no es un recurso abstracto e ilimitado & La interfaz de entrega de la energía y la del cómputo son distintas \\
Gilded Age & La infraestructura de red trae concentración del mercado & El acceso y el interés público requieren atención institucional & La estructura regulatoria global actual es distinta \\
\bottomrule
\end{tabular}
\caption{El uso correcto de la analogía histórica: transferir mecanismos sin perder de vista las diferencias físicas e institucionales.}
\end{table}

\subsection{Resumen de la sección}

Los ciclos históricos dejan cinco recordatorios: la innovación de proceso reduce el costo unitario, el rebote de la demanda amplía el consumo total, el retraso en la construcción genera ciclos de precios, el valor de la sobreconstrucción depende de que los activos puedan reutilizarse, y la concentración del mercado impulsa los estándares y la regulación pública. La siguiente sección aplica estos recordatorios a un juicio concreto: el compute actual todavía no puede tratarse como un bien sustituible, como la electricidad.
\section{Por qué la compute aún no es una Fungible Commodity}

En economía y diseño de mercados, \term{fungibility} significa que una unidad de un bien de cierta categoría puede sustituirse por otra unidad, sin que a las partes de la transacción les importe de qué proveedor concreto proviene. La energía eléctrica estandarizada de la red se acerca a esa experiencia; la compute moderna para IA, en cambio, es muy heterogénea: la arquitectura del chip, la capacidad de memoria del acelerador, la interconexión, el tamaño del clúster, la versión del software, la tasa de fallas, la región geográfica y las condiciones contractuales influyen en si una tarea puede ejecutarse. Reducir todos los recursos a ``GPU-hour'' oculta una gran cantidad de elementos no sustituibles.

\subsection{Una GPU no es un kilowatt-hora}

Una tarea de entrenamiento necesita un conjunto de recursos que se cumplan simultáneamente, no chips aislados. Una oferta de cómputo puede simplificarse como un vector de recursos
\[
\mathbf{c}=(n,\,m,\,b,\,l,\,s,\,a),
\]
donde $n$ representa el número de aceleradores que pueden usarse a la vez, $m$ la memoria disponible por tarjeta, $b$ el ancho de banda de interconexión, $l$ la latencia de comunicación, $s$ la compatibilidad de software y controladores, y $a$ la disponibilidad y la capacidad de recuperación ante fallas. Esta oferta solo tiene valor real cuando el vector de requisitos de la tarea se satisface en su totalidad; un valor muy alto en una dimensión no compensa que otra dimensión no cumpla en absoluto.

\begin{figure}[H]
\centering
\includegraphics[width=0.94\textwidth]{images/00-57-34-compute-nonfungible.jpg}
\caption{Las distintas GPU y clústeres difieren notablemente en precio, capacidad y características de disponibilidad; la compute todavía no puede intercambiarse como la electricidad estandarizada.}
\figsource{Video oficial de Stanford Online 00:57:34, `images/00-57-34-compute-nonfungible.jpg`.}
\end{figure}

\begin{knowledgebox}{Lectura de la figura: primero las restricciones de la tarea, después el precio}
La inferencia de un modelo pequeño puede ejecutarse en muchos tipos de GPU, pero el entrenamiento de un modelo grande exige una gran cantidad de tarjetas de la misma generación, interconexión estable y de alta velocidad, memoria suficiente y disponibilidad continua durante largos periodos. Si los puntos de precio de la figura no están vinculados a estos atributos, no permiten determinar qué oferta es más barata. Para el comprador, la falta de sustituibilidad eleva el costo de migración; para el vendedor, las especificaciones heterogéneas dificultan un mercado unificado, la compensación y la planeación de capacidad.
\end{knowledgebox}

\begin{warningbox}{Aviso de primer uso: la HBM y la memoria principal común no son el mismo recurso}
HBM es High Bandwidth Memory, memoria de alto ancho de banda, que normalmente se conecta de forma estrecha con el acelerador mediante encapsulado avanzado y proporciona memoria de alto ancho de banda para el cálculo matricial. La DRAM de servidor, en cambio, funciona sobre todo como memoria principal de la CPU, con capacidad, ancho de banda y rutas de acceso distintos. Aunque dos máquinas indiquen la misma «memoria total», la memoria del acelerador disponible para el entrenamiento y la topología de comunicación pueden ser completamente diferentes.
\end{warningbox}

\subsection{Las cinco condiciones de ingeniería de una commodity}

El ponente descompone un bien sustituible en cinco condiciones: unidad de medida común, interfaz de entrega estándar, interconexión y agrupación, medición, control y liquidación, y sustituibilidad entre proveedores. Estas condiciones no son términos abstractos de la economía: corresponden, respectivamente, a protocolos, planificadores, descripciones de recursos, compromisos de confiabilidad y reglas de mercado. Para que la compute pase de ser un activo escaso y especializado a una infraestructura ampliamente accesible, las cinco capacidades deben madurar juntas.

\begin{figure}[H]
\centering
\includegraphics[width=0.9\textwidth]{images/01-02-00-fungibility.jpg}
\caption{Las cinco condiciones para que un recurso se convierta en una fungible commodity.}
\figsource{Video oficial de Stanford Online 01:02:00, `images/01-02-00-fungibility.jpg`.}
\end{figure}

\begin{importantbox}{Lectura de la figura: a qué capacidad del sistema corresponde cada condición}
La unidad común exige que las especificaciones de los recursos sean comparables; la interfaz de entrega estándar exige que las tareas puedan enviarse a distintos proveedores; interconnection/pooling exige que la capacidad dispersa pueda combinarse; metering, control y settlement exigen registrar con exactitud el uso, aplicar cuotas y completar la liquidación; la sustituibilidad de proveedores exige que la migración no obligue a reescribir todo el software. Si falta cualquiera de ellas, el mercado puede aparentar tener muchas GPU y, en la práctica, no poder entregar esa capacidad a las tareas que la necesitan.
\end{importantbox}

\begin{table}[H]
\centering
\begin{tabular}{L{0.2\textwidth}L{0.3\textwidth}L{0.36\textwidth}}
\toprule
Condición & Implementación en la infraestructura de cómputo & Dificultad actual \\
\midrule
Common unit & Descripción estandarizada de rendimiento, memoria, interconexión y disponibilidad & Una sola cifra de FLOPs o de GPU-hour oculta las diferencias entre tareas \\
Delivery interface & Job spec portable, imágenes, identidad e interfaces de datos & Fragmentación de controladores, compiladores, redes y permisos \\
Interconnection / pooling & Planificación entre nodos, conjuntos de capacidad y migración ante fallas & Restricciones de latencia, región, cumplimiento normativo de datos y topología \\
Metering / control / settlement & Medición confiable, cuotas, SLA, facturación y compensación & Aislamiento multiinquilino, variabilidad del rendimiento y atribución de responsabilidades \\
Supplier substitution & Las cargas de trabajo pueden migrar entre nubes y entre chips & Kernels propietarios, API de servicios y gravedad de los datos \\
\bottomrule
\end{tabular}
\caption{Glosario: convertir la fungibility en interfaces de sistema implementables y verificables.}
\end{table}

\subsection{Por qué los estándares deben surgir junto con las instituciones}

Los estándares técnicos establecen cómo se representan y conectan los recursos; las institutions, por su parte, se encargan de coordinar la adopción, resolver controversias, hacer cumplir los compromisos y proteger el interés público. Aquí, una institution puede ser un organismo de normalización, una alianza industrial, un operador de mercado, una autoridad reguladora o un sistema de compras públicas. Si solo hay documentos sin aplicación, los proveedores dominantes pueden negarse a la interoperabilidad; si solo hay mandatos sin una interfaz viable, las reglas se quedan en el papel.

\begin{figure}[H]
\centering
\includegraphics[width=0.94\textwidth]{images/01-03-00-standards-institutions.jpg}
\caption{En la historia de la infraestructura, la accesibilidad y la estabilidad suelen impulsarse de forma conjunta por estándares técnicos e instituciones de coordinación.}
\figsource{Video oficial de Stanford Online 01:03:00, `images/01-03-00-standards-institutions.jpg`.}
\end{figure}

\begin{knowledgebox}{Lectura de la figura: los estándares resuelven la interfaz; las instituciones, la coordinación}
AC/DC, el despacho de la red eléctrica, TCP/IP, la interconexión telefónica o las reglas de la aviación no son solo formatos técnicos. Requieren pruebas, certificación, resolución de controversias, recuperación de la inversión y mecanismos de responsabilidad pública. La estandarización de la compute también tocará cuestiones como quién publica el rendimiento, quién asume las interrupciones, cómo se mueven los datos entre países y cómo se evita la manipulación del mercado. La tesis central del ponente es que la fungibility no puede depender únicamente de un mejor software de planificación; también requiere instituciones capaces de obligar a los participantes y de mantener reglas comunes.
\end{knowledgebox}

\begin{warningbox}{La agrupación pública no equivale a ignorar la propiedad ni la seguridad}
Incorporar la capacidad ociosa a un conjunto de recursos más grande requiere mecanismos claros de autorización, aislamiento, medición, compensación y salida. El uso compartido obligatorio sin consentimiento destruye los incentivos a la inversión y la seguridad de los datos; la ausencia total de coordinación, en cambio, puede dejar recursos escasos ociosos durante mucho tiempo o acaparados por unos pocos actores. La tarea del diseño institucional es plasmar en reglas aplicables el beneficio público, la confiabilidad del suministro y los derechos de los participantes.
\end{warningbox}

\subsection{Resumen de la sección}

La falta de sustituibilidad de la compute proviene de restricciones multidimensionales de hardware, topología, software y contratos. Convertirla en commodity requiere una unidad común, interfaces estándar, agrupación, medición y liquidación, y capacidad de sustituir proveedores; estas capacidades técnicas, a su vez, necesitan instituciones que impulsen la adopción, hagan cumplir los compromisos y atiendan el interés público. Con esto, la clase pasa del ciclo cerrado de los modelos a la gobernanza de la infraestructura, y también explica por qué el curso invita a ponentes de las capas de chips, nube, modelos, aplicaciones y políticas públicas.
\section{Síntesis y ampliación}

La línea central de esta clase cabe en una sola frase: el crecimiento de las capacidades de la IA no es la curva de un solo modelo, sino una transformación de pila completa impulsada a la vez por context, compute, capital y culture. Context determina con qué tareas y retroalimentación reales entra en contacto el sistema; compute determina la escala de experimentación que puede sostenerse; capital determina con qué rapidez se construye la infraestructura; culture e institutions determinan los objetivos, las responsabilidades y la asignación de recursos. Si se descuida cualquiera de las capas, la expansión de otra puede convertirse en desperdicio o en riesgo.

\begin{figure}[H]
\centering
\includegraphics[width=0.92\textwidth]{images/01-03-45-food-for-thought.jpg}
\caption{Las dos preguntas que el curso deja a los estudiantes: cómo lograr una compute transition pacífica y de qué manera puede participar cada persona.}
\figsource{Video oficial de Stanford Online, 01:03:45, `images/01-03-45-food-for-thought.jpg`.}
\end{figure}

\begin{importantbox}{Lectura de la figura: el cierre no es una pregunta de predicción para espectadores}
La primera pregunta exige desglosar la «transición pacífica» en accesibilidad, interoperabilidad, concentración del mercado, energía, cadenas de suministro transfronterizas, seguridad e interés público; la segunda pide a los estudiantes que propongan acciones que ellos mismos puedan verificar, por ejemplo definir estándares de recursos, construir workloads portátiles, publicar la metodología de un benchmark, investigar la planificación y la liquidación, participar en organismos de estandarización o explicar las consecuencias de una política. El valor del proyecto del curso no está solo en mostrar un modelo, sino en que los estudiantes entren en la conversación real sobre infraestructura.
\end{importantbox}

\begin{knowledgebox}{Nota de clase: considérate un active participant}
Al final, el ponente enfatiza que los estudiantes no necesitan llegar a ser CEO, directores de investigación o reguladores para tener derecho a participar. Describir con claridad los problemas de una interfaz, publicar experimentos reproducibles, documentar migraciones fallidas, proponer borradores de estándares y pedir evidencia a los ponentes de la industria son acciones que pueden cambiar la discusión. La meta de la preparedness no es predecir mejor quién ganará, sino ser capaz de formular mejores preguntas y construir sistemas verificables en un periodo de incertidumbre.
\end{knowledgebox}

\subsection{Ocho conclusiones transferibles}

\begin{enumerate}
  \item Primero traza las dependencias de toda la pila y después discute los avances en una sola capa; los modelos, los productos y la capacidad de cómputo están insertos en un sistema de recursos más amplio.
  \item Trata la scaling law como un modelo empírico: precisa el intervalo observado, los supuestos de extrapolación y las condiciones en que deja de cumplirse.
  \item Desglosa el context en tareas, herramientas, entorno, permisos y retroalimentación; no lo entiendas solo como un prompt largo.
  \item Al auditar RL, revisa al mismo tiempo el environment, la reward, el verifier y los límites de responsabilidad.
  \item Al leer casos de la industria, sigue el punto de entrada, el modelo, el entorno de ejecución, el retorno de datos y el control de la distribución.
  \item En las gráficas macroeconómicas, distingue las señales tecnológicas, de adopción, de capital y de precios; no uses un tipo de evidencia para sostener una conclusión de otro tipo.
  \item Al usar analogías históricas, transfiere los mecanismos, no calendarios definidos; revisa en especial la vida útil de los activos y las condiciones de estandarización.
  \item Para que el compute se convierta en una infraestructura ampliamente accesible, deben madurar juntos los estándares técnicos y las instituciones que los hacen cumplir.
\end{enumerate}

\subsection{Lecturas adicionales}

Para seguir aprendiendo, se sugieren tres rutas reproducibles. Primera: consulta el sitio oficial del curso Stanford CS153 y los videos oficiales de Stanford Online, y compara cómo los invitados posteriores modifican el marco de cuatro restricciones de esta clase. Segunda: elige un coding agent o una tarea de entrenamiento y traza su context flow, su entorno de ejecución, sus señales de verificación y sus autorizaciones de datos. Tercera: compara la job specification, el SLA y los criterios de facturación de dos mercados de nube o de GPU, y comprueba en la práctica si un workload puede migrarse, en lugar de comparar solo los precios de lista.

\begin{knowledgebox}{Recordatorio de las notas: cómo usan esta clase las sesiones posteriores}
Con cada invitado, registra al menos cuatro cosas: la capa del sistema en la que trabaja, el recurso que considera más escaso, el tipo de evidencia que presenta y los problemas que espera que resuelva el mercado o las instituciones. Al terminar las diez semanas, reúne estas respuestas en una misma tabla; los puntos de conflicto suelen enseñar más que los consensos.
\end{knowledgebox}

\subsection{Tarea práctica: auditar una ruta real de Compute}

Para convertir el marco macro de esta clase en una capacidad de sistemas, puedes elegir un servicio de IA que uses realmente y rastrear una solicitud desde su origen hasta los recursos de cómputo. El objetivo de la tarea no es adivinar toda la arquitectura interna del proveedor, sino dejar claro qué hechos son observables, cuáles provienen de documentación pública y cuáles siguen siendo inferencias. El entregable final debe permitir que otra persona revise las interfaces de recursos, el retorno del context y los límites de responsabilidad.

\begin{enumerate}
  \item Traza la ruta de la solicitud desde el cliente, el servicio de la aplicación y la API del modelo hasta el clúster de inferencia, e indica en qué capa se aplican la identidad, los datos y los permisos de las herramientas.
  \item Elige una tarea repetible y registra la latencia, los tipos de falla, la calidad de la salida y el comportamiento de los reintentos; explica si la reward o la condición de éxito puede verificarse automáticamente.
  \item Compara las especificaciones, los precios y los pasos de migración de dos canales de suministro, y revisa si ``GPU-hour'', el token price o el nombre de una API ocultan que algo es insustituible.
  \item Propón una mejora mínima de estándar, por ejemplo una job spec portátil, códigos de falla unificados, divulgación del rendimiento o comprobantes de eliminación de datos, y explica qué organismo impulsaría su adopción.
\end{enumerate}

\begin{importantbox}{Criterios de aceptación}
Una tarea aceptable debe incluir enlaces a las fuentes, pasos experimentales reproducibles, al menos un caso de falla, columnas separadas para hechos e inferencias, y una explicación de su impacto en la seguridad y el interés público. Si la conclusión solo puede formularse como «tal empresa es más fuerte» o «se necesita más capacidad de cómputo», el análisis todavía no ha llegado a las capas de interfaces, mecanismos e instituciones que exige Frontier Systems.
\end{importantbox}

\end{document}
